% chktex-file 35
% chktex-file 21
% chktex-file 3
\documentclass[12pt]{article}
\usepackage[left=2cm,right=2cm,top=2cm,bottom=2cm,bindingoffset=0cm]{geometry}
\usepackage{amssymb}
\usepackage{amsmath}
\usepackage{amsthm}
\usepackage{enumerate}
\usepackage{enumitem}
\usepackage[T1]{fontenc}
\usepackage[utf8]{inputenc}
\usepackage[matrix,arrow,curve]{xy}
\usepackage[colorlinks=true, urlcolor=blue, linkcolor=blue, citecolor=blue,
    pdfborder={0 0 0}]{hyperref}

%-------------------------------------------------------------------------------
\newtheorem{theorem}{Theorem}[section]
\newtheorem{lemma}[theorem]{Lemma}
\newtheorem{proposition}[theorem]{Proposition}
\newtheorem{remark}[theorem]{Remark}
\newtheorem{corollary}[theorem]{Corollary}
\newtheorem{definition}[theorem]{Definition}
\newtheorem{example}[theorem]{Example}

\newcommand{\projtens}{\mathbin{\widehat{\otimes}}}
\newcommand{\convol}{\ast}
\newcommand{\projmodtens}[1]{\mathbin{\widehat{\otimes}}_{#1}}
\newcommand{\isom}[1]{\mathop{\mathbin{\cong}}\limits_{#1}}
%-------------------------------------------------------------------------------

\begin{document}

\begin{center}
    \Large \textbf{$\mathcal{K}(\ell_p(\Lambda))$ is relatively injective as a 
    right $\mathcal{B}(\ell_p(\Lambda))$-module}\\[0.5cm]
    \small {N. T. Nemesh}\\[0.5cm]
\end{center}

\thispagestyle{empty}

\medskip
\textbf{Abstract:} In this note we prove a somewhat unexpected result: the 
Banach space of compact operators on a reflexive $\ell_p$-space is relatively 
injective as a right Banach module over the algebra of bounded linear operators 
on the same space.
\medskip

\textbf{Keywords:} injective Banach module, compact operators, $\ell_p$-space.

\bigskip

%-------------------------------------------------------------------------------
%	Introduction
%-------------------------------------------------------------------------------



\section{Introduction}\label{SectionIntroduction}

Extension problems have been an important topic of functional analysis since its
inception. Such a problem asks whether a morphism defined on a subobject can be 
extended to the whole object; the objects into which every morphism extends are 
precisely the injective ones. Injectivity is therefore the categorical 
counterpart of the solvability of extension problems. For instance, the field of 
complex numbers is an injective object in the category of Banach spaces. All known injective Banach 
spaces are isomorphic to the space of continuous functions on some compact
space~\cite{BlascIvorConstrInjSpCK}. The property of the object in a given 
category to be injective depends on the ``size'' of the category and 
the ``abundance'' of morphisms. For example, if one restricts the category of 
Banach spaces to separable ones, the space $c_0$ of vanishing sequences becomes 
injective~\cite{SobProjmOnc0}. Moreover, Zippin showed 
in~\cite{ZipSepExtProbm} that all injective spaces in the category of separable 
Banach spaces are isomorphic to $c_0$. Similarly, if one restricts the set of 
morphisms to module maps, without changing the class of objects, then for any 
set $\Lambda$ the space $c_0(\Lambda)$ is relatively injective as a 
right $\ell_\infty(\Lambda)$-module~\cite{NemRelInjC0SmodC0S}. These 
results rely on specific Banach-geometric properties of $c_0$-spaces. One might 
therefore expect that the non-commutative counterpart of these results fails: 
that the space of compact operators on a Hilbert space $\mathcal{K}(H)$ is not 
injective, its geometry being quite different. Nevertheless the result holds, 
not only for Hilbert spaces, but for all reflexive $\ell_p$-spaces.

%-------------------------------------------------------------------------------
%	Preliminaries
%-------------------------------------------------------------------------------

\section{Notation and definitions}\label{SectionNotationAndDefinitions} 
Before we proceed to the main topic, we fix some notation and recall a few 
definitions.

If $M$ is a subset of a set $N$, then $\chi_M$ denotes the indicator function
of $M$. If $f:N\to L$ is an arbitrary function, then $f|_M$ denotes its
restriction to $M$.

For given Banach spaces $E$ and $F$, by $\mathcal{B}(E, F)$ we denote the 
space of all bounded linear operators from $E$ to $F$, and 
by $\mathcal{K}(E, F)$ the space of compact operators 
from $E$ to $F$. We use the 
abbreviations $\mathcal{B}(E):=\mathcal{B}(E, E)$ 
and $\mathcal{K}(E):=\mathcal{K}(E, E)$. The symbol $1_E$ stands for the 
identity operator on $E$. By
$
    \langle\cdot,\cdot\rangle: E\times E^*\to\mathbb{C}, (x,f)\mapsto f(x)
$
we denote the standard duality between $E$ and $E^*$. For given $x\in F$ 
and $f\in E^*$, by $x\bigcirc f$ we denote the rank-one operator from $E$ 
to $F$ defined by the equality $(x\bigcirc f)(y)=\langle y, f\rangle x$. One
can easily check that $\Vert x\bigcirc f\Vert=\Vert x\Vert\Vert f\Vert$.

For an infinite set $\Lambda$ we write $c_0(\Lambda)$ for the space of all 
functions $x:\Lambda\to\mathbb{C}$ such that for every $\varepsilon>0$ the 
set $\{\lambda\in\Lambda:|x(\lambda)|\geq\varepsilon\}$ is finite, and 
$\ell_\infty(\Lambda)$ for the space of all bounded such functions. For 
$p\in[1,+\infty)$ the space $\ell_p(\Lambda)$ consists of all 
functions $x:\Lambda\to\mathbb{C}$ with finite $\ell_p$-norm
$\Vert x\Vert_p:=(\sum_{\lambda\in\Lambda}|x(\lambda)|^p)^{1/p}$. For a 
family of Banach spaces $\{E_\lambda:\lambda\in\Lambda\}$ and $p\in[1,+\infty]$
by $\bigoplus_p\{E_\lambda:\lambda\in\Lambda\}$ we denote 
their $\ell_p$-sum. Clearly, 
$\ell_p(\Lambda)=\bigoplus_p\{\mathbb{C}:\lambda\in\Lambda\}$.

A bounded linear functional $\mu\in\ell_\infty(\Lambda)^*$ is called 
a \textit{singular state} on $\ell_\infty(\Lambda)$ if (i) $\Vert\mu\Vert=1$; 
(ii) $\mu(\chi_\Lambda)=1$; and (iii) $\mu|_{c_0(\Lambda)}=0$. Conditions (i) 
and (ii) say that $\mu$ is a state on the unital Banach algebra 
$\ell_\infty(\Lambda)$, while (iii) says that $\mu$ vanishes on $c_0(\Lambda)$. 
Singular states exist for every infinite $\Lambda$: the quotient 
$\ell_\infty(\Lambda)/c_0(\Lambda)$ is a unital commutative $C^*$-algebra, hence 
carries a state $\sigma$ (for instance the evaluation at any point of its 
maximal ideal space), and the composition of $\sigma$ with the quotient 
homomorphism $\ell_\infty(\Lambda)\to\ell_\infty(\Lambda)/c_0(\Lambda)$ is a 
singular state on $\ell_\infty(\Lambda)$. We do not require $\mu$ to be 
multiplicative.

Now we shall recall a few notions from the geometry of Banach spaces. A 
sequence $(x_n)_{n\in\mathbb{N}}$ in a Banach space $E$ is called 
\textit{weakly null}, \textit{weak* null} or \textit{norm null} if it converges 
to zero weakly, weakly* or in norm, respectively. The 
sequence $(x_n)_{n\in\mathbb{N}}$ is called \textit{semi-normalized} if the norms 
of its elements have a positive lower bound and a finite upper bound. 
Let $(\Lambda, <)$ be a well-ordered set. A family of 
vectors $(e_\lambda)_{\lambda\in\Lambda}$ is called a \textit{Schauder basis} 
of $E$ if for any vector $x\in E$ there exists a unique family of scalars 
$(f_\lambda(x))_{\lambda\in\Lambda}\subset\mathbb{C}$ such that 
\[
    x=\sum_{\lambda\in\Lambda}f_\lambda(x) e_\lambda.
\]
In general the value of this sum may depend on the ordering of $\Lambda$; if it 
does not, the basis is called \textit{unconditional}. One can 
show [\cite{KalAlbTopicsBanSpTh}, theorem 1.1.3] that the 
maps $(f_\lambda)_{\lambda\in\Lambda}$ are well-defined bounded linear 
functionals. They are called the \textit{biorthogonal functionals}
of $(e_\lambda)_{\lambda\in\Lambda}$. 

A sequence of 
vectors $(e_n)_{n\in\mathbb{N}}$ in $E$ is called a \textit{basic sequence} if 
it is a Schauder basis of its closed linear span. The following construction 
gives an important source of basic sequences. 
Suppose $(e_\lambda)_{\lambda\in\Lambda}$ is a Schauder basis of some Banach 
space $E$, and let $(\Lambda_n)_{n\in\mathbb{N}}$ be a sequence of disjoint 
finite subsets of $\Lambda$ such that $\max(\Lambda_n)<\min(\Lambda_{n+1})$ for 
all $n\in\mathbb{N}$. A sequence of non-zero vectors $(u_n)_{n\in\mathbb{N}}$ 
such that $u_n$ lies in the linear span of $(e_\lambda)_{\lambda\in\Lambda_n}$ 
is called a \textit{block-basic sequence}. Such sequences are indeed 
basic [\cite{KalAlbTopicsBanSpTh}, lemma 1.3.5]. If $E$ and $F$ are two 
Banach spaces, then two sequences $(u_n)_{n\in\mathbb{N}}\subset E$ 
and $(v_n)_{n\in\mathbb{N}}\subset F$ are called \textit{equivalent} if there 
is an isomorphism of Banach spaces $T:E\to F$ such that $T(u_n)=v_n$ for 
all $n\in\mathbb{N}$.

Finally, we give a short introduction to Banach homology theory. Let $A$ 
be a Banach algebra. We shall consider only right Banach modules over
$A$ with contractive outer action $\cdot:X\times A\to X$. If $X$ and $Y$ are two
right Banach $A$-modules, then a map $\phi:X\to Y$ is an \textit{$A$-morphism}
if it is a continuous $A$-module map. Banach $A$-modules and $A$-morphisms form
a category which we denote by $\mathbf{mod}-A$.

In $\mathbf{mod}-A$ the notion of injectivity can be defined in different ways. 
We shall be concerned with relative injectivity. Let $\xi:X\to Y$ be an 
$A$-morphism. 
Then it is called $c$-\textit{relatively admissible} if $T\circ \xi=1_X$ 
for some bounded linear operator $T:Y\to X$ with $\Vert T\Vert\leq c$. 
A Banach $A$-module $J$ is called $C$-\textit{relatively injective} if for any 
$c$-relatively admissible $A$-morphism $\xi:X\to Y$ and 
any $A$-morphism $\phi:X\to J$ there exists a continuous $A$-module 
map $\psi:Y\to J$ such that $\Vert\psi\Vert\leq cC\Vert\phi\Vert$ and the 
diagram
\[
    \xymatrix{
    & {Y} \\
    {J} \ar@{<--}[ur]^{\psi} &{X} \ar[u]^{\xi} \ar[l]^{\phi}}  % chktex 3
\]
commutes. A right Banach $A$-module is called \textit{relatively injective} 
if it is $C$-relatively injective for some $C\geq 1$. As an example, 
if $A=\{0\}$, the category $\mathbf{mod}-A$ reduces to the ordinary category of 
Banach spaces, and in this case all Banach spaces are relatively injective. 

All injectivity statements in this note are verified by means of a single 
criterion, due to Helemskii, which reduces relative injectivity to the 
existence of one module retraction. Suppose $A$ is a unital Banach algebra 
acting on itself by multiplication, and let $J$ be a right Banach $A$-module. 
The space $\mathcal{B}(A, J)$ of bounded linear operators from $A$ to $J$ is a 
right $A$-module under the action $(\mathcal{T}\cdot a)(b):=\mathcal{T}(ab)$; it 
is the \textit{cofree} module cogenerated by $J$. There is a canonical 
$A$-morphism
\[
    \rho_J: J\to\mathcal{B}(A, J),\quad x\mapsto(a\mapsto x\cdot a),
\]
and $\rho_J$ always admits a bounded linear (not necessarily module) left 
inverse, namely $\mathcal{T}\mapsto\mathcal{T}(1_A)$, so it is relatively 
admissible. Helemskii's criterion [\cite{HelBanLocConvAlg}, proposition 
VII.1.39(III)] states that $J$ is relatively injective if and only if 
$\rho_J$ admits a left inverse \textit{in the category} $\mathbf{mod}-A$, that 
is, a right $A$-module map $\tau:\mathcal{B}(A, J)\to J$ 
with $\tau\circ\rho_J=1_J$. Moreover, by a theorem of 
White [\cite{WhiteInjmoduAlg}, proposition 3.8], if such a $\tau$ can be chosen 
contractive then $J$ is $1$-relatively injective.

%-------------------------------------------------------------------------------
%   Relative injectivity of the right $\mathcal{B}(H)$-module $\mathcal{K}(H)$
%-------------------------------------------------------------------------------

\section{Main result}\label{SecionMainResult}

In this section we shall prove that for any reflexive $\ell_p$-space the 
right Banach $\mathcal{B}(\ell_p(\Lambda))$-module $\mathcal{K}(\ell_p(\Lambda))$ 
is relatively injective.

Before turning to compact operators, we recall that the ambient algebra is 
always relatively injective as a module over itself, and explain why the 
passage to the ideal of compact operators is harder.

\begin{proposition}\label{BEModBEIsRelInj}
    For any Banach space $E$ the right $\mathcal{B}(E)$-module 
    $\mathcal{B}(E)$ is $1$-relatively injective.
\end{proposition}
\begin{proof} If $E=\{0\}$ then $\mathcal{B}(E)=\{0\}$ and the claim is trivial, 
    so assume $E\neq\{0\}$. Fix $e_0\in E$ and $f_0\in E^*$ with 
    $\langle e_0, f_0\rangle = 1$ and $\Vert e_0\Vert=\Vert f_0\Vert=1$, and 
    define
    \[
        \tau:
        \mathcal{B}(\mathcal{B}(E), \mathcal{B}(E))\to\mathcal{B}(E),\quad
        \mathcal{T}\mapsto(x\mapsto \mathcal{T}(x\bigcirc f_0)(e_0)).
    \]
    Using the identity $A\circ(x\bigcirc f_0)=(Ax)\bigcirc f_0$ together with 
    $\langle e_0, f_0\rangle=1$, it is trivial to check that $\tau$ is a 
    contractive right $\mathcal{B}(E)$-module map which is a left inverse for the 
    canonical morphism $\rho_{\mathcal{B}(E)}$. By Helemskii's criterion and 
    White's theorem, recalled in 
    Section~\ref{SectionNotationAndDefinitions}, the module $\mathcal{B}(E)$ 
    is $1$-relatively injective.
\end{proof}

The single evaluation works because $\tau(\mathcal{T})$ may be any bounded 
operator. This fails in the compact case. If 
$\mathcal{T}:\mathcal{B}(E)\to\mathcal{K}(E)$ and we keep the same formula, then 
$\tau(\mathcal{T})$ need not be compact: for $\mathcal{T}(S):=(Se_0)\bigcirc f_0$ 
the formula gives $\tau(\mathcal{T})=1_E$, which is non-compact in infinite 
dimensions. To keep the value in $\mathcal{K}(E)$ we average the evaluation over 
a whole biorthogonal system instead of reading it off at one point. The singular 
state below provides this average.

There is also an a priori obstruction that explains why compactness cannot be 
taken for granted. We record it as a necessary condition for the ideal of 
compact operators.

\begin{proposition}\label{KNotComplemented}
    Let $E$ be an infinite-dimensional Banach space. If the right 
    $\mathcal{B}(E)$-module $\mathcal{K}(E)$ is relatively injective, then 
    $\mathcal{K}(E)$ is not complemented in $\mathcal{B}(E)$ as a Banach space.
\end{proposition}
\begin{proof} The ideal $\mathcal{K}(E)$ has no left identity: if $u\in
    \mathcal{K}(E)$ satisfied $u\circ K=K$ for all $K\in\mathcal{K}(E)$, then 
    testing on the rank-one operators $K=x\bigcirc f$ would give $u(x)=x$ for 
    all $x\in E$, that is $u=1_E$, which is not compact in infinite 
    dimensions. On the other hand, by a result of 
    Ramsden [\cite{RamsdenThesis}, proposition 2.2.8(iii)], a relatively 
    injective closed right ideal of a Banach algebra which is complemented as a 
    subspace possesses a left identity. Were $\mathcal{K}(E)$ relatively 
    injective and complemented in $\mathcal{B}(E)$, this would furnish a left 
    identity for $\mathcal{K}(E)$, a contradiction.
\end{proof}

Proposition~\ref{KNotComplemented} is a necessary condition only, and a weak 
one, since $\mathcal{K}(E)$ is uncomplemented in $\mathcal{B}(E)$ for most 
classical $E$. It is not vacuous, however. If the Calkin algebra 
$\mathcal{B}(E)/\mathcal{K}(E)$ is finite-dimensional, then $\mathcal{K}(E)$ is 
finite-codimensional and hence complemented, so by 
Proposition~\ref{KNotComplemented} its module $\mathcal{K}(E)$ is not relatively 
injective. Such spaces exist: on the Argyros-Haydon space $\mathfrak{X}_{AH}$ 
every operator is a scalar plus a compact one, so the Calkin algebra is 
one-dimensional \cite{ArgyrosHaydon}; Tarbard constructed preduals of $\ell_1$ 
whose Calkin algebra is finite-dimensional of any prescribed dimension 
\cite{Tarbard}. The point of the proposition is that 
relative injectivity of $\mathcal{K}(E)$ cannot come from a bounded projection 
and must be witnessed by a module splitting of the cofree module. We construct 
such a splitting for reflexive $\ell_p$-spaces in the rest of this section.

\begin{proposition}\label{LInftyStateByMeanProps} Let $\Lambda$ be an infinite 
    set, and $\mu$ be a singular state on $\ell_\infty(\Lambda)$. 
    If $|\mu(x)|\geq c$ for 
    some $x\in\ell_\infty(\Lambda)$ and $c>0$, then for any $c'<c$ the 
    set $\{\lambda\in\Lambda: |x(\lambda)|>c'\}$ is infinite.
\end{proposition}
\begin{proof} Suppose to the contrary that for some $c'<c$ the 
    set $F:=\{\lambda\in\Lambda: |x(\lambda)|>c'\}$ is finite. Decompose 
    $x=x\chi_F+x\chi_{\Lambda\setminus F}$. Since $F$ is finite, 
    $x\chi_F\in c_0(\Lambda)$, hence $\mu(x\chi_F)=0$ by property $(iii)$ of a 
    singular state. On the complement we 
    have $\Vert x\chi_{\Lambda\setminus F}\Vert\leq c'$, so 
    \[
        |\mu(x)|
        =|\mu(x\chi_F)+\mu(x\chi_{\Lambda\setminus F})|
        =|\mu(x\chi_{\Lambda\setminus F})|
        \leq\Vert\mu\Vert\Vert x\chi_{\Lambda\setminus F}\Vert
        \leq c' < c,
    \]
    which contradicts $|\mu(x)|\geq c$. Therefore $F$ is infinite.
\end{proof}

\begin{proposition}\label{EInftyStateByMeanProps} Let $E$ be a reflexive
    Banach space, let $\Lambda$ be an infinite set, and let $\mu$ be a singular 
    state on $\ell_\infty(\Lambda)$. Then there exists a contractive linear 
    operator 
    $
        M_{\mu}:
        \bigoplus\nolimits_\infty\{ E : \lambda\in\Lambda \}\to E
    $
    such that: 
    \begin{enumerate}[label = (\roman*)]
        \item 
            $\langle M_{\mu}(\xi), f\rangle
            =\mu\left(
                \bigoplus_\infty\{ \langle \xi(\lambda), f\rangle: 
                    \lambda\in\Lambda\}
            \right)
            $ for all $\xi\in \bigoplus_\infty\{ E : \lambda\in\Lambda \}$ 
            and $f \in E^*$;
        \item 
            $
            M_{\mu}\left(
                \bigoplus_\infty\{ x: \lambda\in\Lambda\}
            \right) = x
            $ for any $x \in E$;
    \end{enumerate} 
\end{proposition}
\begin{proof} Let $\xi\in \bigoplus_\infty\{ E : \lambda\in\Lambda \}$, 
    then $\{\langle \xi(\lambda), f\rangle:\lambda\in\Lambda\}\subset\mathbb{C}$ 
    is bounded, so the map
    \[
        F_{\mu,\xi}: E^*\to\mathbb{C},\, 
        f\mapsto \mu\left(
            \bigoplus\nolimits_\infty\{\langle \xi(\lambda), f\rangle: 
                \lambda\in\Lambda\}
        \right)
    \]
    is well defined. Clearly, it is a bounded linear functional 
    with $\Vert F_{\mu,\xi}\Vert\leq\Vert \xi\Vert$. Since $E$ is 
    reflexive, there exists a unique vector $M_{\mu,\xi}\in E$ such 
    that $\Vert M_{\mu,\xi}\Vert=\Vert F_{\mu,\xi}\Vert$ 
    and $F_{\mu,\xi}(f)=\langle M_{\mu,\xi}, f\rangle$ 
    for all $f\in E^*$. Consider the map
    $
        M_{\mu}:
        \bigoplus\nolimits_\infty\{ E : \lambda\in\Lambda \}\to E,\, 
        \xi\mapsto M_{\mu,\xi}.
    $
    It is easy to check that $M_{\mu}$ is a linear operator. It is 
    contractive because 
    $
        \Vert M_{\mu}(\xi)\Vert
        =\Vert F_{\mu,\xi}\Vert
        \leq\Vert \xi\Vert.
    $
    We now verify the two properties of $M_{\mu}$.

    $(i)$ This equality follows directly from the defining property of the 
    vector $M_{\mu,\xi}$.

    $(ii)$ Using the equality $\mu(\chi_\Lambda)=1$, for 
    any $x\in E$ and $f\in E^*$ we have 
    \[
        \left\langle M_{\mu}\left(
            \bigoplus\nolimits_\infty\{ x: \lambda\in\Lambda\}
            \right),
            f
        \right\rangle
        =\mu\left(\bigoplus\nolimits_\infty\left\{
            \left\langle x, f\right\rangle:\lambda\in\Lambda
        \right\}\right)
        =\langle x,f \rangle \mu(\chi_{\Lambda})
        =\langle x,f \rangle. 
    \]
    Since $f\in E^*$ is arbitrary, we obtain the desired equality.
\end{proof}

\begin{proposition}\label{TUTMapIsLinAndBnd}
    Let $E$ be a reflexive Banach space. 
    Let ${(e_\lambda)}_{\lambda\in\Lambda}\subset E$ 
    and ${(f_\lambda)}_{\lambda\in\Lambda}\subset E^*$ be two families such
    that $\sup\{\Vert e_\lambda\Vert\Vert f_\lambda\Vert: \lambda\in\Lambda\}$ 
    is finite. Then for any singular state $\mu$ on $\ell_\infty(\Lambda)$ and 
    any bounded linear operator $\mathcal{T}:\mathcal{B}(E)\to\mathcal{B}(E)$ the 
    map
    \[
        T_{\mu,\mathcal{T}}: 
        E\to E,\, 
        x\mapsto M_{\mu}\left(
            \bigoplus\nolimits_{\infty}\{ 
                \mathcal{T}(x\bigcirc f_\lambda)(e_\lambda) : \lambda\in\Lambda 
            \}
        \right)
    \]
    is a well-defined bounded linear operator with norm at 
    most 
    $
        \Vert \mathcal{T} \Vert
        \sup\{ \Vert e_\lambda\Vert \Vert f_\lambda\Vert : \lambda\in\Lambda\}.
    $
\end{proposition}
\begin{proof}
    For $x\in E$, denote 
    $
        \xi_{\mathcal{T},x}=\bigoplus\nolimits_{\infty}\{ 
            \mathcal{T}(x\bigcirc f_\lambda)(e_\lambda) : \lambda\in\Lambda 
        \}.
    $ 
    Clearly, 
    \[
        \sup\{\Vert \xi_{\mathcal{T},x}(\lambda)\Vert:\lambda\in\Lambda\}
        \leq
        \Vert \mathcal{T}\Vert \Vert x\Vert
        \sup\{ \Vert e_\lambda\Vert \Vert f_\lambda\Vert : \lambda\in\Lambda\},
    \]
    so 
    $
        \xi_{\mathcal{T},x}\in
            \bigoplus\nolimits_{\infty}{}\{E:\lambda\in\Lambda\}
    $ 
    and the value $M_{\mu}(\xi_{\mathcal{T},x})$ is well defined; hence so is 
    the map $T_{\mu,\mathcal{T}}$. It is easy to check 
    that $T_{\mu,\mathcal{T}}$ is linear. Furthermore, using 
    proposition~\ref{EInftyStateByMeanProps}, for any $x\in E$
    we get
    \[
        \Vert T_{\mu,\mathcal{T}}(x)\Vert
        \leq\Vert M_{\mu}\Vert\Vert \xi_{\mathcal{T},x}\Vert
        \leq\Vert \mathcal{T}\Vert \Vert x\Vert
        \sup\{ \Vert e_\lambda\Vert \Vert f_\lambda\Vert : \lambda\in\Lambda\}.
    \] 
    This proves the norm estimate for $T_{\mu,\mathcal{T}}$.
\end{proof}

\begin{proposition}\label{BlckBasForNonCompOp}
    Let $E$ be a reflexive Banach space with a Schauder 
    basis $(e_\lambda)_{\lambda\in\Lambda}$. Let $T:E\to E$ be a non-compact 
    bounded linear operator. Then there exists a weakly null, semi-normalized 
    block-basic 
    sequence $(u_n)_{n\in\mathbb{N}}$ of $(e_\lambda)_{\lambda\in\Lambda}$
    such that $(T(u_n))_{n\in\mathbb{N}}$ is semi-normalized.
\end{proposition} 
\begin{proof} Note that the property of a sequence to be weakly null, norm null
    or semi-normalized is preserved when passing to subsequences or 
    equivalent sequences. We shall use these facts many times during the proof.
    
    Since $T$ is not compact, there exists a weakly null 
    sequence $(x_n)_{n\in\mathbb{N}}\subset E$ such 
    that $(T(x_n))_{n\in\mathbb{N}}$ is not norm null. Passing to a subsequence,
    we may assume that $\inf\{\Vert T(x_n)\Vert: n\in\mathbb{N}\}>0$. Note 
    that $(x_n)_{n\in\mathbb{N}}$ cannot be norm null, since 
    otherwise $(T(x_n))_{n\in\mathbb{N}}$ would be norm null. Therefore, passing
    to a subsequence once again, we may additionally assume  
    that $\inf\{\Vert x_n\Vert: n\in\mathbb{N}\}>0$. As $(x_n)_{n\in\mathbb{N}}$
    is weakly convergent, it is bounded; so is $(T(x_n))_{n\in\mathbb{N}}$, 
    because $T$ is bounded. Thus we have constructed a weakly null, 
    semi-normalized sequence $(x_n)_{n\in\mathbb{N}}$ such 
    that $(T(x_n))_{n\in\mathbb{N}}$ is also semi-normalized.

    Since $(x_n)_{n\in\mathbb{N}}$ is a countable family, it lies in the 
    closed linear span of some countable 
    subfamily $(e_\lambda)_{\lambda\in\Lambda'}$ of the 
    basis $(e_\lambda)_{\lambda\in\Lambda}$. By the Bessaga--Pe{\l}czy\'nski 
    selection principle [\cite{KalAlbTopicsBanSpTh}, proposition 1.3.10] there 
    exists a block-basic sequence $(u_k)_{k\in\mathbb{N}}$ 
    of $(e_\lambda)_{\lambda\in\Lambda'}$ which is equivalent to some 
    subsequence of $(x_n)_{n\in\mathbb{N}}$. As $(u_k)_{k\in\mathbb{N}}$ is a 
    block-basic sequence of $(e_\lambda)_{\lambda\in\Lambda'}$, it is also a 
    block-basic sequence of the larger 
    basis $(e_\lambda)_{\lambda\in\Lambda}$. Recall 
    that $(x_n)_{n\in\mathbb{N}}$ is weakly null and semi-normalized, 
    with $(T(x_n))_{n\in\mathbb{N}}$ also semi-normalized. The 
    sequence $(u_k)_{k\in\mathbb{N}}$ inherits these properties, since it is 
    equivalent to a subsequence of $(x_n)_{n\in\mathbb{N}}$.
\end{proof}

\begin{proposition}\label{TUTMapIsComp}
    Let $\Lambda$ be a set, let $p\in(1,+\infty)$, and 
    denote $E=\ell_p(\Lambda)$. Let $(e_\lambda)_{\lambda\in\Lambda}\subset E$ 
    be the standard basis and $(f_\lambda)_{\lambda\in\Lambda}\subset E^*$ the 
    respective biorthogonal functionals. Then for any singular state $\mu$ 
    on $\ell_\infty(\Lambda)$ and any bounded linear 
    operator $\mathcal{T}:\mathcal{B}(E)\to\mathcal{K}(E)$ the 
    map $T_{\mu, \mathcal{T}}$ is a compact linear operator of norm at 
    most $\Vert \mathcal{T}\Vert$. 
\end{proposition}
\begin{proof} If $E$ is finite-dimensional, then 
    $\mathcal{B}(E)=\mathcal{K}(E)$, so there is nothing to prove.

    Suppose $E$ is infinite-dimensional, so that $\Lambda$ is infinite, and 
    denote $T:=T_{\mu, \mathcal{T}}$. By 
    proposition~\ref{TUTMapIsLinAndBnd}, the map $T$ is a bounded 
    linear operator with $\Vert T\Vert\leq\Vert \mathcal{T}\Vert$.
    Assume, towards a contradiction, that $T$ is not compact. 
    By proposition~\ref{BlckBasForNonCompOp} we obtain a block-basic 
    sequence $(u_n)_{n\in\mathbb{N}}$ of $(e_\lambda)_{\lambda\in\Lambda}$ such
    that $c:=\inf\{\Vert T(u_n)\Vert:n\in\mathbb{N}\}>0$ 
    and $C:=\sup\{\Vert u_n\Vert:n\in\mathbb{N}\}<+\infty$.

    Fix $n\in\mathbb{N}$. By the Hahn-Banach theorem there exists $g_n\in E^*$ 
    such that $\Vert g_n\Vert=1$ and 
    $
        \langle T(u_n), g_n\rangle=\Vert T(u_n)\Vert
    $.
    Consider the vector 
    $
        s_n:=\bigoplus\nolimits_\infty\{
            \langle \mathcal{T}(u_n\bigcirc f_\lambda)(e_\lambda), g_n\rangle: 
            \lambda\in\Lambda
        \}.
    $
    Note that $s_n\in\ell_\infty(\Lambda)$, 
    because 
    $
        \sup\{|s_n(\lambda)| : \lambda\in\Lambda\} \leq C\Vert \mathcal{T}\Vert.
    $ 
    It is easy to check that
    $
        \mu(s_n)=\langle T(u_n), g_n\rangle
    $,
    so $|\mu(s_n)|=\mu(s_n)\geq c$. By 
    proposition~\ref{LInftyStateByMeanProps}, the 
    set $A_n:=\{\lambda\in\Lambda: |s_n(\lambda)|> c/2\}$ is infinite.
    Since the sets $A_n$ are infinite for all $n\in\mathbb{N}$, by induction we 
    can construct a sequence of distinct 
    elements $(\lambda_n)_{n\in\mathbb{N}}\subset\Lambda$ 
    with $\lambda_n\in A_n$ for all $n\in\mathbb{N}$. In other words, there 
    exists a sequence of distinct 
    elements $(\lambda_n)_{n\in\mathbb{N}}\subset\Lambda$ such that
    $
        |\langle 
            \mathcal{T}(u_n\bigcirc f_{\lambda_n})(e_{\lambda_n}), g_n
        \rangle| > c/2.
    $

    Since $(u_n)_{n\in\mathbb{N}}$ is a block basis of the standard basis 
    of $E$, the supports of these vectors are disjoint, so for 
    any $x\in\ell_\infty(\mathbb{N})$ and $y\in\ell_p(\Lambda)$ we get
    \[
        \begin{aligned}
            \sum_{n\in\mathbb{N}}\Vert x(n) \langle y, f_{\lambda_n}\rangle u_n \Vert
            =\left(
                \sum_{n\in\mathbb{N}}
                    |x(n)|^p|\langle y, f_{\lambda_n}\rangle|^p \Vert u_n \Vert^p
            \right)^{1/p}
            \leq C \Vert x\Vert \left(
                \sum_{n\in\mathbb{N}}|\langle y, f_{\lambda_n}\rangle|^p
            \right)^{1/p}
            \leq C \Vert x\Vert\Vert y\Vert
        \end{aligned}.
    \]
    This means that the series 
    $
        \sum_{n\in\mathbb{N}} x(n) \langle y, f_{\lambda_n}\rangle u_n
    $
    is absolutely convergent, and a fortiori convergent.
    Therefore, for each $x\in\ell_\infty(\mathbb{N})$ we obtain a well-defined 
    map 
    \[
        S_x: 
        E\to E, \,
        y\mapsto \sum_{n\in\mathbb{N}} x(n)(u_n\bigcirc f_{\lambda_n})(y).
    \]
    One can easily check that $S_x$ is a linear operator, and the estimate above 
    shows that $\Vert S_x\Vert\leq C\Vert x\Vert$. Now for each $n\in\mathbb{N}$ 
    we define a map
    \[
        f_n: 
        \ell_\infty(\mathbb{N})\to\mathbb{C}, \,
        x\mapsto \langle \mathcal{T}(S_x)(e_{\lambda_n}), g_n\rangle.
    \]
    Clearly, it is a bounded linear functional. The 
    sequence $(e_{\lambda_n})_{n\in\mathbb{N}}$ is weakly null, and by definition 
    of $\mathcal{T}$ the operator $\mathcal{T}(S_x)$ is 
    compact, so the sequence 
    $
        (\mathcal{T}(S_x)(e_{\lambda_n}))_{n\in\mathbb{N}}
    $ is norm null. As the functionals $(g_n)_{n\in\mathbb{N}}\subset E^*$ are 
    norm bounded, we conclude that
    \[
        \lim_{n\to\infty}f_n(x)
        =\lim_{n\to\infty}\langle \mathcal{T}(S_x)(e_{\lambda_n}), g_n\rangle
        =0.
    \]
    This means that $(f_n)_{n\in\mathbb{N}}$ is a weak* null sequence. On the 
    other hand, 
    \[
        \sum_{k\in\mathbb{N}}|f_n(\chi_{\{k\}})|
        \geq |f_n(\chi_{\{n\}})|
        =|\langle 
            \mathcal{T}(u_n\bigcirc f_{\lambda_n})(e_{\lambda_n}), g_n
        \rangle|
        > c/2.
    \]
    So $\liminf_{n\to\infty}\sum_{k\in\mathbb{N}}|f_n(\chi_{\{k\}})|>0$,
    which contradicts Phillips' lemma [\cite{PhilOnLinTransfrm}, corollary 3.4].
    Therefore $T$ is compact.
\end{proof}

\begin{proposition}\label{TauMapProps}
    Let $\Lambda$ be a set, let $p\in(1,+\infty)$, and 
    denote $E=\ell_p(\Lambda)$. Let $(e_\lambda)_{\lambda\in\Lambda}\subset E$ 
    be the standard basis and $(f_\lambda)_{\lambda\in\Lambda}\subset E^*$ the 
    respective biorthogonal functionals. Let $\mu$ be any 
    singular state on $\ell_\infty(\Lambda)$. Define
    \[
        \tau_{\mu}:
        \mathcal{B}(\mathcal{B}(E), \mathcal{K}(E))\to\mathcal{K}(E), \,
        \mathcal{T}\mapsto T_{\mu, \mathcal{T}}.
    \]
    Then $\tau_{\mu}$ is a contractive $\mathcal{B}(E)$-morphism of 
    right $\mathcal{B}(E)$-modules such 
    that $\tau_{\mu}\circ\rho_{\mathcal{K}(E)}=1_{\mathcal{K}(E)}$, 
    where 
    $
        \rho_{\mathcal{K}(E)}:
        \mathcal{K}(E)\to\mathcal{B}(\mathcal{B}(E), \mathcal{K}(E)), \,
        K\mapsto(T\mapsto K\circ T)
    $.
\end{proposition}
\begin{proof} By proposition~\ref{TUTMapIsComp}, the map $\tau_{\mu}$ is 
    well defined. It is routine to check that $\tau_{\mu}$ is linear. 
    From the same proposition we also obtain 
    that $\Vert \tau_{\mu}(\mathcal{T})\Vert\leq\Vert \mathcal{T}\Vert$, 
    so $\tau_{\mu}$ is contractive. For 
    any $\mathcal{T}\in\mathcal{B}(\mathcal{B}(E), \mathcal{K}(E))$, 
    $A\in\mathcal{B}(E)$ and $x\in E$ we have
    \[
        \begin{aligned}
            \tau_{\mu}(\mathcal{T}\cdot A)(x)
            &=M_{\mu}\left(
                \bigoplus\nolimits_{\infty}\{ 
                    (\mathcal{T}\cdot A)(x\bigcirc f_\lambda)(e_\lambda) : 
                    \lambda\in\Lambda 
                \}
            \right)\\
            &=M_{\mu}\left(
                \bigoplus\nolimits_{\infty}\{ 
                    \mathcal{T}(A(x)\bigcirc f_\lambda)(e_\lambda) : 
                    \lambda\in\Lambda 
                \}
            \right)\\
            &=\tau_{\mu}(\mathcal{T})(A (x))\\
            &=(\tau_{\mu}(\mathcal{T})\cdot A)(x).
        \end{aligned}
    \]
    Therefore $\tau_{\mu}$ is a right $\mathcal{B}(E)$-module map. 
    Finally, using proposition~\ref{EInftyStateByMeanProps} and the 
    biorthogonality of the families $(e_\lambda)_{\lambda\in\Lambda}$ 
    and $(f_\lambda)_{\lambda\in\Lambda}$, we obtain, for 
    all $K\in\mathcal{K}(E)$ and $x\in E$,
        \[
        \begin{aligned}
            \tau_{\mu}(\rho_{\mathcal{K}(E)}(K))(x)
            &=M_{\mu}\left(
                \bigoplus\nolimits_{\infty}\{ 
                    \rho_{\mathcal{K}(E)}(K)(x\bigcirc f_\lambda)(e_\lambda) : 
                    \lambda\in\Lambda 
                \}
            \right)\\
            &=M_{\mu}\left(
                \bigoplus\nolimits_{\infty}\{ 
                    \langle e_\lambda, f_\lambda\rangle K(x)  : 
                    \lambda\in\Lambda 
                \}
            \right)\\
            &=M_{\mu}\left(
                \bigoplus\nolimits_{\infty}\{ 
                     K(x): \lambda\in\Lambda 
                \}
            \right)\\
            &=K(x).
        \end{aligned}
    \]
    In other words, $\tau_{\mu}\circ\rho_{\mathcal{K}(E)}=1_{\mathcal{K}(E)}$.
\end{proof}

\begin{theorem}\label{BHModKHIsRelInj}
    Let $E$ be an $\ell_p$-space for some $p\in(1,+\infty)$. Then the 
    right $\mathcal{B}(E)$-module $\mathcal{K}(E)$ is $1$-relatively injective.
\end{theorem}
\begin{proof} By assumption $E=\ell_p(\Lambda)$ for some set $\Lambda$. If 
    $\Lambda$ is finite, then $E$ is finite-dimensional, 
    so $\mathcal{K}(E)=\mathcal{B}(E)$, and this module is $1$-relatively 
    injective by Proposition~\ref{BEModBEIsRelInj}. 
    Assume now that $\Lambda$ is infinite, and fix a singular state $\mu$ 
    on $\ell_\infty(\Lambda)$, which exists by the discussion in 
    Section~\ref{SectionNotationAndDefinitions}. By 
    proposition~\ref{TauMapProps} there exists a contractive right 
    $\mathcal{B}(E)$-morphism $\tau_{\mu}$ which is a left
    inverse for 
    $
        \rho_{\mathcal{K}(E)}:
        \mathcal{K}(E)\to\mathcal{B}(\mathcal{B}(E), \mathcal{K}(E)), \,
        K\mapsto(T\mapsto K\circ T).
    $
    By Helemskii's criterion and White's theorem 
    (Section~\ref{SectionNotationAndDefinitions}), the contractivity 
    of $\tau_{\mu}$ shows that the right $\mathcal{B}(E)$-module 
    $\mathcal{K}(E)$ is $1$-relatively injective.
\end{proof}


%-------------------------------------------------------------------------------
%   Bibliography
%-------------------------------------------------------------------------------

\begin{thebibliography}{999}
    %
    \bibitem{BlascIvorConstrInjSpCK}
    \textit{Blasco, J.L. and Ivorra, C.}, On constructing injective spaces of
    type $C(K)$, Indagationes Mathematicae, 9(12), 161--172, 1998.
    %
    %
    \bibitem{SobProjmOnc0}
    \textit{Sobczyk, A.}, Projection of the space ($m$) on its subspace ($c_0$),
    Bulletin of the American Mathematical Society, 47(12), 938--947, 1941.
    %
    %
    \bibitem{ZipSepExtProbm}
    \textit{Zippin, M.}, The separable extension problem, Israel Journal of
    Mathematics, 26(3--4), 372--387, 1977.
    %
    %
    \bibitem{NemRelInjC0SmodC0S}
    \textit{Nemesh, N.T.}, A note on relatively injective $C_0(S)$-modules, 
    Functional Analysis and Its Applications, 55, 298--303, 2021.
    %
    %
    \bibitem{PhilOnLinTransfrm}
    \textit{Phillips, R.S.}, On linear transformations,
    Transactions of American Mathematical Society, 48(3), 516--541, 1940.
    %
    %
    \bibitem{HelBanLocConvAlg}
    \textit{A.Ya. Helemskii}, Banach and locally convex algebras. Oxford
    University Press, (1993)
    %
    %
    \bibitem{KalAlbTopicsBanSpTh}
    \textit{Albiac F., Kalton, N.J.} Topics in Banach space theory, 
    Springer, 2006.
    %
    %
    \bibitem{WhiteInjmoduAlg}\textit{M. C. White}, Injective modules for uniform 
    algebras, Proceedings of London Mathematical Society, 3(1), p. 155--184, 
    1991.
    %
    %
    \bibitem{RamsdenThesis}\textit{P. Ramsden}, Homological properties of 
    semigroup algebras, PhD thesis, University of Leeds, 2009.
    %
    %
    \bibitem{ArgyrosHaydon}\textit{S. A. Argyros and R. G. Haydon}, A 
    hereditarily indecomposable $\mathcal{L}_\infty$-space that solves the 
    scalar-plus-compact problem, Acta Mathematica, 206(1), 1--54, 2011. 
    \url{https://doi.org/10.1007/s11511-011-0058-y}
    %
    %
    \bibitem{Tarbard}\textit{M. Tarbard}, Hereditarily indecomposable, separable 
    $\mathcal{L}_\infty$ Banach spaces with $\ell_1$ dual having few but not very 
    few operators, Journal of the London Mathematical Society, 85(3), 737--764, 
    2012. \url{https://arxiv.org/abs/1011.4776}
    %

\end{thebibliography}

Norbert Nemesh, Faculty of Mechanics and Mathematics, Moscow State University,
Moscow 119991 Russia

\textit{E-mail:} nemeshnorbert@yandex.ru


\end{document}  % chktex 17